% Part: proof-theory
% Chapter: natural-deduction
% Section: rules-proofs

\documentclass[../../../include/open-logic-section]{subfiles}

\begin{document}

\olfileid{pt}{nat}{rp}

\olsection{Aturan lan \usetoken{P}{proof}}

Kita miwiti kanthi menehi sawetara dhéfinisi lan netepaké sawetara istilah.

Coba delengen rong aturan ing \Log{N1c} lan~\Log{N1i} kang nyebut $!A
\lif !B$.
\[
\bottomAlignProof
\AxiomC{$\Discharge{!A}{x}$}
\shortDeduce
\DeduceC{$!B$}
\DischargeRule{\Intro{\lif}}{x}
\UnaryInfC{$!A \lif !B$}
\DisplayProof
\qquad
\bottomAlignProof
\AxiomC{$!A \lif !B$}
\AxiomC{$!A$}
\RightLabel{\Elim{\lif}}
\BinaryInfC{$!B$}
\DisplayProof
\]
Aturan \Elim{\lif} cukup prasaja. !!{formula}s ing ndhuwur garis
iku \emph{premis}. !!{formula} ing sisih kiwa iku
!!{formula} \emph{utama}~$!A \lif !B$. Kabèh aturan !!{elimination}
nduwèni siji premis kaya mangkono, tansah ing sisih paling kiwa.
Premis iku diarani premis \emph{mayor} saka inferènsi. Yèn ana
premis liya, kaya ing kéné, premis iku diarani premis \emph{minor}.
Premis aturan !!{introduction} uga diarani premis mayor.
!!{formula} ing sangisoré garis iku \emph{konklusi}.

Aturan \Intro{\lif} mung nduwèni siji premis, lan ing kéné konklusiné
yaiku formula utama~$!A \lif !B$. Kabèh aturan !!{introduction}
kaya mangkono: konklusiné dadi formula utama, lan operator pokoké
yaiku operator kang ``dilebokaké.''

Notasi $\Discharge{!A}{x}$ nduwèni teges kaya mangkéné. Ing
!!a{proof}, inferènsi \Intro{\lif} ditrapaké marang !!a{formula}~$!B$
minangka premis. Sub-!!{proof} kang pungkasané~$!B$ nduwèni sawetara
!!{formula}s wiwitan. Iki diarani \emph{asumsi}. !!^a{proof}
kanggo~$!B$ saka siji utawa luwih asumsi kang wujudé~$!A$ iku
!!a{proof} yèn $!A$ ndhèrèkaké~$!B$. Nalika aturan \Intro{\lif}
ditrapaké, kita nyimpulaké~$!A \lif !B$, lan konklusi iku ora gumantung
manèh marang asumsi~$!A$. Kanggo nandhani iki, asumsi kang wujudé~$!A$
ing !!{proof} kang ngemot inferènsi~\Intro{\lif} mau \emph{dibusak}
utawa \emph{ditutup}. Aturané nemtokaké asumsi endi kang bisa dibusak.
Ing \Intro{\lif} kanthi konklusi $!A \lif !B$, asumsi iku
wujudé~$!A$, yaiku padha karo antesèdené implikasi. Kadhang kala
perlu dituduhaké asumsi endi kang dibusak déning aturan endi;
kanggo kuwi dienggo label pambusakan~$x$. Sing wigati, aturan kang
mbusak asumsi ora \emph{kudu} mbusak \emph{kabèh} asumsi kang wujudé
cocog, malah ora kudu mbusak asumsi siji-sijia. Tuladhané, iki bener
sanajan \Intro{\lif} kapisan ora mbusak asumsi apa-apa:
\[
\AxiomC{$\Discharge{!A}{y}$}
\DischargeRule{\Intro{\lif}}{x}
\UnaryInfC{$!B \lif !A$}
\DischargeRule{\Intro{\lif}}{y}
\UnaryInfC{$!A \lif (!B \lif !A)$}
\DisplayProof
\]
Yèn sawijining aturan ora mbusak asumsi apa-apa, label pambusakané
ora ditulis (kaya ing kéné).

Ing !!{proof} iki
\[
\AxiomC{$\Discharge{!A}{x}$}
\AxiomC{$\Discharge{!A}{y}$}
\RightLabel{\Intro{\land}}
\BinaryInfC{$!A \land !A$}
\RightLabel{\Elim{\land}}
\UnaryInfC{$!A$}
\DischargeRule{\Intro{\lif}}{x}
\UnaryInfC{$!A \lif !A$}
\DischargeRule{\Intro{\lif}}{y}
\UnaryInfC{$!A \lif (!A \lif !A)$}
\DisplayProof
\]
inferènsi \Intro{\lif} kapisan mbusak asumsi sisih kiwa
$!A^x$, déné kang kapindho mbusak asumsi sisih tengen~$!A^y$.
Tegesé, kabèh asumsi mawa label~$x$ dibusak déning \Intro{\lif}
mawa label~$x$, lan kabèh asumsi mawa label~$y$ déning aturan
mawa label~$y$. Supaya iki lumaku, kabèh asumsi kang labelé padha
uga kudu dadi !!{formula} kang padha.

Aturan kuantor rada luwih ruwet. Tuladhané, aturan ing
\Log{N1i} kanggo $\lexists$ yaiku:
\[
\bottomAlignProof
\AxiomC{$!A(t)$}
\RightLabel{\Intro{\lexists}}
\UnaryInfC{$\lexists[x][!A(x)]$}
\DisplayProof
\qquad
\bottomAlignProof
\AxiomC{$\lexists[x][!A(x)]$} 
\AxiomC{$\Discharge{!A(c)}{x}$}
\shortDeduce
\DeduceC{$!B$}
\DischargeRule{\Elim{\lexists}}{x}
\BinaryInfC{$!B$} \DisplayProof 
\bottomAlignProof
\] Ing aturan
\Intro{\lexists}, premisé~$!A(t)$ kanthi $t$ minangka suku katutup,
yaiku suku tanpa variabel. Inferènsi kaya mangkono bener---tegesé,
cocog karo skéma aturan \Intro{\lexists}---yèn ana !!{formula}~$!A$,
variabel~$x$, lan suku katutup~$t$ nganti konklusiné
$\lexists[x][!A]$ lan premisé~$\Subst{!A}{t}{x}$.
Aturan \Elim{\lexists} ngidini kita mbusak asumsi kang wujudé~$!A(c)$:
ing saben inferènsi ana !!{constant}~$c$, lan asumsi kang wujudé
$\Subst{!A}{c}{x}$ bisa dibusak. Kabèh asumsi kang dibusak déning
siji inferènsi kudu nganggo~$c$ kang padha. Konstanta iki diarani
\emph{variabel eigen} saka inferènsi \Elim{\lexists}. Inferènsi iku
mung bener yèn $c$~ora katon ing asumsi kabuka liyané kajaba
$!A(c)$ kang dibusak, ora katon ing~$!B$, lan ora katon ing~$!A$.
Iki diarani \emph{syarat variabel eigen}.

!!^{proof}s ing \Log{N1c} lan~\Log{N1i} iku wit winates saka
formula kang kawangun sacara induktif saka aksioma nganggo
aturan inferènsi. Saben pucuk wit iku asumsi, kang bisa mawa label
pambusakan; saben !!{formula} liyané ndhèrèk saka !!{formula}s
ing sadhuwuré langsung lumantar salah siji aturan inferènsi ing
sistem mau. Yèn sawijining kemunculan !!{formula} minangka asumsi
ing wit sadhuwuré inferènsi durung dibusak déning inferènsi ing
antarané, asumsi iku diarani \emph{kabuka} (ing inferènsi mau).

\begin{defn}\ollabel{defn:proof}
\emph{!!^{proof}s} ing \Log{N1c} lan~\Log{N1i} iku wit winates saka
!!{formula}s kang ditetepaké sacara induktif kaya mangkéné:
\begin{enumerate}
  \item\ollabel{defn:prf:assumption} Saben !!{formula}
  kang mawa label pambusakan kayata $x$ utawa $y$, dhéwé wis dadi
  !!a{proof}. Ing !!a{proof} kang mung isi siji !!{formula},
  !!{formula} iku dianggep asumsi kabuka.
  \item\ollabel{defn:prf:one-prem} Yèn $R$ aturan kang nduwèni siji,
  loro, utawa telung premis $!A_1$, $!A_2$, $!A_3$, kanthi
  konklusi~$!A$, lan $\delta_1$, $\delta_2$, $\delta_3$ iku
  !!{proof}s kang pungkasané $!A_1$, $!A_2$, $!A_3$,
  miturut urutané, banjur
  \[
  \AxiomC{}
  \RightLabel{$\delta_1$}
  \DeduceC{$!A_1$}
  \RightLabel{$R$}
  \UnaryInfC{$!A$}
  \DisplayProof
\qquad
  \AxiomC{}
  \RightLabel{$\delta_1$}
  \DeduceC{$!A_1$}
  \AxiomC{}
  \RightLabel{$\delta_2$}
  \DeduceC{$!A_2$}
  \RightLabel{$R$}
  \BinaryInfC{$!A$}
  \DisplayProof
\qquad
  \AxiomC{}
  \RightLabel{$\delta_1$}
  \DeduceC{$!A_1$}
  \AxiomC{}
  \RightLabel{$\delta_2$}
  \DeduceC{$!A_2$}
  \AxiomC{}
  \RightLabel{$\delta_3$}
  \DeduceC{$!A_3$}
  \RightLabel{$R$}
  \TrinaryInfC{$!A$}
  \DisplayProof
  \]
  saben-sabené dadi !!a{proof}, angger label asumsi ing !!{proof}s mau
  cocog (ora ana rong !!{formula}s kang béda mawa label pambusakan
  kang padha) lan inferènsiné ngetrapaké~$R$ kanthi bener.
  \item Kemunculan !!{formula} paling dhuwur ing !!{proof} anyar
  dianggep \emph{asumsi kabuka} saka !!{proof} iku yèn uga dadi
  asumsi kabuka saka $\delta_1$, $\delta_2$, utawa~$\delta_3$,
  kajaba kemunculan kang dibusak déning aturan mau. Yèn aturané
  mawa label~$x$ (utawa $x\, y$), kanggo \Intro{\lif}
  lan~\Intro{\lnot} kang dibusak yaiku kabèh asumsi kabuka
  ing~$\delta_1$ mawa label~$x$; kanggo \Elim{\lexists} yaiku
  kabèh asumsi kabuka mawa label~$x$ ing~$\delta_2$; kanggo
  \Elim{\lor} yaiku kabèh asumsi kabuka mawa label~$x$
  ing~$\delta_2$ lan kabèh asumsi kabuka mawa label~$y$
  ing~$\delta_3$.
\end{enumerate}
\end{defn}

!!{proof}s kang dienggo ing pambangunan induktif sawijining
!!a{proof} diarani \emph{sub-!!{proof}s}-é. Sub-!!{proof}s kang
pungkasané ana ing premis-premis sawijining inferènsi diarani
\emph{sub-!!{proof}s langsung} saka inferènsi mau. Aturan
\Log{N1c} lan \Log{N1i} katon ing \olref{tab:N1}.

\subfile{rules-N1}

\begin{defn}
\emph{!!{height}}~$\pheight{\delta}$ saka !!a{proof}~$\delta$
ditetepaké sacara induktif kaya mangkéné.
\begin{enumerate}
  \item Yèn $\delta$ mung isi asumsi, $\pheight{\delta} = 0$.
  \item Yèn $\delta$ pungkasané inferènsi kanthi siji premis lan
  sub-!!{proof} langsung~$\delta_1$, banjur $\pheight{\delta} = \pheight{\delta_1} + 1$.
  \item Yèn $\delta$ pungkasané inferènsi kanthi rong premis lan
  sub-!!{proof}s langsung~$\delta_1$ lan~$\delta_2$, banjur $\pheight{\delta} =
  \max(\pheight{\delta_1}, \pheight{\delta_2}) + 1$.
  \item Yèn $\delta$ pungkasané \Elim{\lor} kanthi
  sub-!!{proof}s langsung~$\delta_1$, $\delta_2$, lan~$\delta_3$, banjur $\pheight{\delta} =
  \max(\pheight{\delta_1}, \pheight{\delta_2}, \pheight{\delta_3}) + 1$.
\end{enumerate}
Yèn sekuen~$S$ nduwèni !!a{proof} kang dhuwuré $\le n$,
kita nulis $\Proves[n] S$.
\end{defn}

\begin{ex}
Ing \Log{N1i}, saben !!{formula} bisa dadi !!a{proof} saka awaké dhéwé
minangka asumsi kabuka. Mula
\[
!A \land !B^x
\]
iku !!a{proof}~$\delta_1$. Dhuwuré~$0$.
Saka $\delta_1$ kita bisa ngasilaké rong !!{proof}s~$\delta_2$
lan~$\delta_3$ kanthi ngetrapaké salah siji saka rong versi~\Elim{\land}
\[
\AxiomC{$!A \land !B^x$}
\RightLabel{\Elim{\land}[2]}
\UnaryInfC{$!B$}
\DisplayProof
\qquad
\AxiomC{$!A \land !B^x$}
\RightLabel{\Elim{\land}[1]}
\UnaryInfC{$!A$}
\DisplayProof
\]
Sub-!!{proof}s langsung saka inferènsi pungkasan ing~$\delta_2$
lan~$\delta_3$ iku padha, yaiku $!A \land !B$.
Kaloro !!{proof}s mau nduwèni kemunculan $!A\land !B$
minangka asumsi kabuka. Kita olèh
$\pheight{\delta_2} = \pheight{\delta_3} = 1$.

Aturan \Intro{\land} mbutuhaké rong premis kang wujudé $!B$ lan $!A$
kanggo mbeneraké konklusi kang wujudé~$!B \land !A$. Mula, ing
ngisor iki !!a{proof}~$\delta_4$:
\[
\AxiomC{$!A \land !B^x$}
\RightLabel{\Elim{\land}[2]}
\UnaryInfC{$!B$}
\LeftSubproofLabel{$\delta_2$}
\AxiomC{$!A \land !B^x$}
\RightLabel{\Elim{\land}[1]}
\UnaryInfC{$!A$}
\RightSubproofLabel{$\delta_3$}
\RightLabel{\Intro{\land}}
\BinaryInfC{$!B \land !A$}
\DisplayProof
\]
Dhuwuré $\pheight{\delta_4} = \max(\pheight{\delta_2},
\pheight{\delta_3})+1 = \max(1,1) + 1 = 2$. Ana rong kemunculan
$!A \land !B$ minangka asumsi, lan kaloroné kabuka.

Pungkasan, kita bisa ngetrapaké \Intro{\lif} kanggo ngasilaké $\delta_5$:
\[
\AxiomC{$\Discharge{!A \land !B}{x}$}
\RightLabel{\Elim{\land}[2]}
\UnaryInfC{$!B$}
\LeftSubproofLabel{$\delta_2$}
\AxiomC{$\Discharge{!A \land !B}{x}$}
\RightLabel{\Elim{\land}[1]}
\UnaryInfC{$!A$}
\RightSubproofLabel{$\delta_3$}
\RightLabel{\Intro{\land}}
\BinaryInfC{$!B \land !A$}
\DischargeRule{\Intro{\lif}}{x}
\UnaryInfC{$(!A \land !B) \lif (!B \land !A)$}
\DisplayProof
\]
Dhuwuré $\pheight{\delta_4} + 1 = 3$. Inferènsi \Intro{\lif}
mbusak kabèh kemunculan asumsi kang wujudé $!A \land !B$
mawa label~$x$, yaiku kaloro asumsi kang sadurungé kabuka
ing subbukti langsung. Amarga iku mung siji-sijiné jinis
asumsi ing $\delta_5$, saiki ora ana asumsi kabuka, mula iki
!!a{proof} \emph{kanggo} !!{formula} pungkasané
$(!A \land !B) \lif (!B \land !A)$.
\end{ex}

\end{document}
